% =============================================================================
%  Pipeline Walkthrough — a reverse-engineered account of the implementation
%
%  Every number in this document was produced by instrumenting the actual code
%  in this repository and recording the counts at each step.
% =============================================================================
\documentclass[11pt,a4paper]{article}

\usepackage[T1]{fontenc}
\usepackage[utf8]{inputenc}
\usepackage[margin=2.2cm]{geometry}
\usepackage{booktabs}
\usepackage{longtable}
\usepackage{array}
\usepackage{xcolor}
\usepackage{graphicx}
\usepackage{enumitem}
\usepackage{caption}
\usepackage{amsmath}
\usepackage{amssymb}
\usepackage{fancyvrb}
\usepackage[hidelinks]{hyperref}

\definecolor{codebg}{HTML}{F4F6F8}
\definecolor{warn}{HTML}{B7770B}

\newcommand{\code}[1]{\texttt{\small #1}}
\newcommand{\arr}{$\rightarrow$}

\setlist[itemize]{itemsep=2pt,topsep=4pt}
\setlist[enumerate]{itemsep=2pt,topsep=4pt}
\captionsetup{font=small,labelfont=bf}

\DefineVerbatimEnvironment{codeblock}{Verbatim}
  {fontsize=\small,frame=leftline,framerule=1pt,rulecolor=\color{gray},xleftmargin=8pt,framesep=6pt}

\newenvironment{keypoint}
  {\begin{center}\begin{minipage}{0.95\textwidth}%
   \vspace{4pt}\hrule\vspace{7pt}\small}
  {\vspace{7pt}\hrule\vspace{4pt}\end{minipage}\end{center}}

\title{\textbf{Pipeline Walkthrough}\\[0.4em]
       \large Probabilistic Motor Subtype Progression Model for Parkinson's Disease\\[0.6em]
       \normalsize A step-by-step account of how every result was generated}
\author{\small Prepared as a study and defence guide}
\date{August 2026}

\begin{document}
\maketitle

\begin{abstract}
\noindent
This document reverse-engineers the implementation in this repository and
traces every reported number back to the code that produced it. Nothing here is
generic methodology: each count was obtained by instrumenting the actual
functions and recording the data shape before and after every transformation.

Where a design decision has a documented rationale, that rationale is quoted.
Where it does not, this is stated explicitly rather than filled in with a
plausible-sounding reason.
\end{abstract}

\tableofcontents
\clearpage

% =============================================================================
\section{Current project state}
% =============================================================================

\subsection{What exists}

Twenty-four tracked files: seven modules in \code{src/ppmi\_subtype/}
(1{,}552 lines), three scripts, three test files, one configuration file, and
documentation.

\begin{table}[h]
\centering
\begin{tabular}{@{}lrl@{}}
\toprule
\textbf{Location} & \textbf{Lines} & \textbf{Role} \\
\midrule
\code{src/ppmi\_subtype/} & 1{,}552 & All analysis logic \\
\code{scripts/}           &   596 & Orchestration: call \code{src}, print, save \\
\code{tests/}             &   598 & 74 tests, all exercising \code{src} \\
\bottomrule
\end{tabular}
\caption{Code inventory. The scripts contain no analysis logic.}
\end{table}

\subsection{What has actually been executed}

Established by comparing file modification times, not assumed:

\begin{table}[h]
\centering
\begin{tabular}{@{}lr@{}}
\toprule
\textbf{Item} & \textbf{Last modified} \\
\midrule
Newest source file (\code{scripts/00\_data\_audit.py}) & 2026-08-21 11:15 \\
All 20 output files                                    & 2026-08-21 11:17 \\
\bottomrule
\end{tabular}
\end{table}

\noindent
Every output post-dates every source file, so \textbf{no output is stale}. All
74 tests pass against the same code that produced those outputs.

\subsection{Which script produced which output}

\begin{table}[h]
\centering
\small
\begin{tabular}{@{}lll@{}}
\toprule
\textbf{Output file} & \textbf{Produced by} & \textbf{Line} \\
\midrule
\code{outputs/visit\_panel.csv}          & \code{01\_build\_panel.py} & 118 \\
\code{outputs/baseline.csv}              & \code{01\_build\_panel.py} & 119 \\
\code{outputs/panel\_summary.txt}        & \code{01\_build\_panel.py} & 94 \\
\code{outputs/validation.txt}            & \code{01\_build\_panel.py} & 112 \\
\code{outputs/transitions/*.csv} (9)     & \code{02\_transition\_analysis.py} & 113--222 \\
\code{outputs/transitions/report.txt}    & \code{02\_transition\_analysis.py} & 270 \\
\code{outputs/transitions/figures/*.png} (6) & \code{plots.py} via stage 2 & 59 \\
\bottomrule
\end{tabular}
\end{table}

\subsection{Dead code and unused configuration}

Three findings, stated plainly because a reviewer will find them:

\begin{enumerate}
  \item \code{config.yaml} line 19 defines \code{participant\_status:}. A grep
        across \code{src/}, \code{scripts/} and \code{tests/} returns zero
        references. It is a configured path with no consumer.
  \item Of the seventeen files under \code{data/}, exactly \textbf{four} are
        opened by the pipeline. The remainder --- demographics, cohort history,
        five dictionary files, and four of the six files in
        \code{data/validation/} --- are human reference material that no code
        path reads.
  \item \code{scripts/00\_data\_audit.py} contains no \code{to\_csv} or
        \code{write\_text} call. It writes nothing. It is a read-only
        diagnostic, \emph{not} a pipeline stage.
\end{enumerate}

\noindent
\textbf{Not established from the current repository:} there are no logs, no
notebooks, no shell scripts and no run manifests. Execution history is
inferable only from file timestamps and git commits.

% =============================================================================
\section{Execution pipeline}
% =============================================================================

\subsection{The order is proven, not assumed}

The numeric filename prefixes suggest an order, but the order is
\emph{established} by a data dependency:
\code{scripts/02\_transition\_analysis.py} line 85 reads
\code{outputs/visit\_panel.csv}, and the only code that writes that file is
\code{scripts/01\_build\_panel.py} line 118.

\begin{codeblock}
        +-----------------------------------------+
        |  00_data_audit.py    (DIAGNOSTIC ONLY)  |
        |  reads raw data, writes nothing         |
        +-----------------------------------------+
                (optional; off the critical path)

  RAW DATA --> 01_build_panel.py --> outputs/visit_panel.csv
                                              |
                                              v
                               02_transition_analysis.py
                                              |
                                              v
                                    outputs/transitions/
\end{codeblock}

\subsection{Module dependency graph}

Derived from the \code{import} statements; the graph is acyclic.

\begin{codeblock}
  config.py        <- foundation; imports nothing internal
      ^
      +-- loaders.py      (config)
      +-- subtype.py      (config)
      +-- validate.py     (config, loaders)
      +-- panel.py        (config, loaders, subtype)

  transitions.py   <- imports NOTHING internal. Pure DataFrame functions.
  plots.py         <- imports NOTHING internal. Pure DataFrame functions.
\end{codeblock}

\noindent
\code{transitions.py} and \code{plots.py} do not read the configuration at all.
They take DataFrames in and return DataFrames or figures out. That is why they
carry the heaviest test coverage: they have no hidden state.

% =============================================================================
\section{Stage 0 --- the raw inputs}
% =============================================================================

\subsection{The four files the pipeline opens}

\subsubsection*{File 1 --- the visit skeleton}

\code{data/clinical/PPMI\_Curated\_Data\_Cut\_Public\_20260511.xlsx}, sheet
\code{20260511}.

\begin{table}[h]
\centering
\begin{tabular}{@{}lr@{}}
\toprule
Shape & \textbf{19{,}450 rows $\times$ 212 columns} \\
Unique \code{PATNO} & 4{,}788 \\
Distinct \code{EVENT\_ID} values & 18 \\
Duplicate \code{(PATNO, EVENT\_ID)} & \textbf{0} \\
\bottomrule
\end{tabular}
\end{table}

\noindent
\textbf{One row = one patient at one scheduled study visit.}

\begin{codeblock}
 PATNO EVENT_ID  YEAR visit_date    subgroup  COHORT  age_at_visit  td_pigd  pigd   hy
100001       BL     0    10/2020 Sporadic PD       1     67.386301      1.0   0.4  2.0
100001      V04     1    11/2021 Sporadic PD       1     68.512329      1.0   0.4  2.0
100001      V06     2    11/2022 Sporadic PD       1     69.528767      1.0   0.4  2.0
\end{codeblock}

Read by \code{loaders.load\_curated\_cut()} (lines 45--48).

\textbf{Columns used:} \code{PATNO}, \code{EVENT\_ID} (merge keys),
\code{subgroup} (cohort filter), \code{YEAR} and \code{visit\_date} (time
axis), \code{td\_pigd} (validation reference only), plus the baseline predictor
columns written to \code{baseline.csv}.

\begin{keypoint}
This file already contains PPMI's own derived scores (\code{pigd}, \code{hy},
\code{updrs3\_score}, \code{td\_pigd}). \textbf{The pipeline does not use any of
them to build the label.} It uses this file purely as a visit skeleton plus
covariates, and computes the subtype independently from raw items. That
independence is what makes the validation in Section~\ref{sec:validation}
meaningful --- otherwise it would be comparing a number against itself.
\end{keypoint}

\subsubsection*{File 2 --- patient questionnaire items}

\code{data/motor\_updrs/MDS\_UPDRS\_Part\_II\_\_Patient\_Questionnaire\_20Aug2026.csv}

\begin{table}[h]
\centering
\begin{tabular}{@{}lr@{}}
\toprule
Shape & \textbf{37{,}777 rows $\times$ 22 columns} \\
Unique \code{PATNO} & 5{,}341 \\
Duplicate \code{(PATNO, EVENT\_ID)} & \textbf{0} \\
\bottomrule
\end{tabular}
\end{table}

\noindent
\textbf{One row = one patient's self-reported questionnaire at one visit.}

All 22 columns: \code{REC\_ID, PATNO, EVENT\_ID, PAG\_NAME, INFODT, NUPSOURC,
NP2SPCH, NP2SALV, NP2SWAL, NP2EAT, NP2DRES, NP2HYGN, NP2HWRT, NP2HOBB,
NP2TURN, NP2TRMR, NP2RISE, NP2WALK, NP2FREZ, NP2PTOT, ORIG\_ENTRY,
LAST\_UPDATE}.

\textbf{The pipeline reads exactly three item columns:} \code{NP2TRMR}
(tremor), \code{NP2WALK} (walking and balance), \code{NP2FREZ} (freezing). The
other ten items are not used. Read by \code{loaders.load\_updrs\_part2()}.

\subsubsection*{File 3 --- motor examination items}

\code{data/motor\_updrs/MDS-UPDRS\_Part\_III\_20Aug2026.csv}

\begin{table}[h]
\centering
\begin{tabular}{@{}lr@{}}
\toprule
Shape & \textbf{39{,}310 rows $\times$ 65 columns} \\
Unique \code{PATNO} & 5{,}344 \\
Duplicate \code{(PATNO, EVENT\_ID)} & \textcolor{warn}{\textbf{6{,}907}} \\
\bottomrule
\end{tabular}
\end{table}

\noindent
\textbf{One row = one clinician-administered motor examination --- not one row
per visit.}

\begin{keypoint}
\textbf{The most important structural fact in the raw data.} Part II has zero
duplicate \code{(PATNO, EVENT\_ID)} pairs; Part III has 6{,}907. The reason is
in the columns: \code{PDSTATE} (\code{ON}/\code{OFF}), \code{PDTRTMNT},
\code{ONEXAM}, \code{OFFEXAM}, \code{ONOFFORDER}. A medicated patient can be
examined \emph{twice at the same visit} --- once off medication and once on it
--- producing two rows with different scores:

\begin{Verbatim}[fontsize=\small,xleftmargin=10pt]
PATNO   EVENT_ID  PDSTATE  PDTRTMNT  NP3GAIT  NP3PSTBL
3002    V06       OFF      1         2        1          <- exam A
3002    V06       ON       1         1        0          <- exam B, same visit
\end{Verbatim}

Merging this against the visit skeleton without handling it would either
duplicate the visit or silently keep whichever row happened to come first. Both
are wrong. This is precisely why \code{select\_medication\_state()} and
\code{deduplicate\_visits()} exist.
\end{keypoint}

\textbf{The pipeline reads thirteen item columns plus three metadata columns:}
\code{NP3PTRMR, NP3PTRML, NP3KTRMR, NP3KTRML, NP3RTARU, NP3RTALU, NP3RTARL,
NP3RTALL, NP3RTALJ, NP3RTCON} (tremor), \code{NP3GAIT, NP3FRZGT, NP3PSTBL}
(gait and postural), plus \code{PDSTATE}, \code{PDTRTMNT}, \code{NHY}. The
remaining $\approx 46$ columns are not used.

\subsubsection*{File 4 --- the validation reference}

\code{data/validation/PPMI\_Original\_Cohort\_BL\_to\_Year\_5\_Dataset\_Apr2020.csv}

\begin{table}[h]
\centering
\begin{tabular}{@{}lr@{}}
\toprule
Shape & \textbf{3{,}445 rows $\times$ 159 columns} \\
Unique \code{PATNO} & 683 \\
\code{td\_pigd\_old} & 1.0: 1{,}440 \quad 2.0: 550 \quad 3.0: 1{,}014 \quad NaN: 441 \\
\bottomrule
\end{tabular}
\end{table}

\noindent
This file is \textbf{never used for analysis}. It is read once, by
\code{validate.against\_archived\_three\_class()}, purely to check the
recomputed labels. It is the only file in the repository carrying a
\textbf{three-class} subtype label, because the current curated cut collapsed
its \code{td\_pigd} to two classes.

\subsection{Why these sixteen items?}

\textbf{What the code does.} \code{config.yaml} lines 42--60 list eleven tremor
items and five PIGD items.

\textbf{Documented rationale}, quoted verbatim from \code{config.yaml} lines
36--39: \emph{``Definition follows Stebbins et al., Movement Disorders 2013, and
matches the derivation notes in the Curated Data Cut dictionary for
\code{td\_pigd} verbatim.''}

\textbf{Interpretation.} The item selection is not a modelling choice this
project made; it is a published clinical definition that was cross-checked
against PPMI's own derivation notes.

\textbf{Assumption introduced.} That the Stebbins definition is the right
phenotype to model. The proposal commits to this explicitly --- its stated goal
is to \emph{retain} the clinically validated ratio rather than invent a new
score.

% =============================================================================
\section{Stage 1 --- \code{01\_build\_panel.py}}
% =============================================================================

\subsection{A. Purpose}

\textbf{Scientific question:} \emph{For every visit of every sporadic PD
patient in PPMI, what was that patient's motor subtype at that moment, and can
we prove the label is correct?}

This script converts raw MDS-UPDRS item scores into a longitudinal, labelled
patient-visit panel, and validates that panel against PPMI's own derivations
before any analysis is permitted to use it.

\subsection{B. Inputs}

Three files (Section~3): the curated cut, MDS-UPDRS Part II, MDS-UPDRS
Part III. A fourth, the archived cohort, is read only for validation.

\subsection{C. Code walkthrough with actual counts}

The script's real work is one call, \code{build\_visit\_panel(cfg)} in
\code{panel.py} line 56, which decomposes as follows.

\subsubsection*{Step 1.1 --- \code{filter\_cohort()} \quad (\code{panel.py}:44)}

\begin{codeblock}
out = frame[frame["subgroup"].isin(keep)]
out = out[~out["EVENT_ID"].isin(set(cohort["excluded_events"]))]
out = out[out["YEAR"] <= cohort["max_follow_up_years"]]
\end{codeblock}

\begin{table}[h]
\centering
\begin{tabular}{@{}lrr@{}}
\toprule
\textbf{Operation} & \textbf{Rows} & \textbf{Removed} \\
\midrule
\code{load\_curated\_cut()}                          & 19{,}450 & --- \\
\code{subgroup} $\in$ \{Sporadic PD\}                & \phantom{0}5{,}789 & 13{,}661 \\
\code{EVENT\_ID} $\notin$ \{ST, PW, U01, U02\}       & \phantom{0}5{,}774 & \phantom{00,0}15 \\
\code{YEAR} $\leq$ 15                                & \phantom{0}5{,}774 & \phantom{00,00}0 \\
\bottomrule
\end{tabular}
\caption{Cohort filtering. Patients: 4{,}788 \arr\ 1{,}080.}
\end{table}

\noindent
\textbf{What happened to the removed rows.} The 13{,}661 rows dropped at the
first step are the other PPMI subgroups: Hyposmia (1{,}759 patients), RBD
(572), LRRK2 (400), Healthy Control (330), GBA (310), SWEDD (60) and others.
They are not deleted from disk; they are simply not selected. The 15 further
rows are visits labelled \code{ST} (unscheduled) or \code{PW} (premature
withdrawal) --- verified: those are the only two excluded codes actually
present. The \code{YEAR} filter removes nothing, because the curated cut
already stops at year 15.

\textbf{Why this step exists.} \code{config.yaml} lines 114--116 document it as
\emph{``Sporadic PD is the primary analysis population; the genetic PD subgroups
are listed so a sensitivity analysis can switch them on without touching
code.''}

\textbf{The scientific rationale for restricting to sporadic PD is not
explicitly documented in the repository.} A defensible interpretation --- label
it as interpretation, not as a recorded decision --- is that genetic PD
subgroups (LRRK2, GBA, SNCA) have known differences in motor phenotype and
progression rate, so pooling them would mix populations. The config is written
so this can be tested rather than assumed.

\subsubsection*{Step 1.2 --- \code{select\_medication\_state()} \quad (\code{subtype.py}:159)}

\begin{codeblock}
state = frame[state_col].astype("string").str.strip().str.upper()
matches = state == cfg.medication_state          # "OFF"
if treat_col in frame.columns:
    untreated = pd.to_numeric(frame[treat_col], errors="coerce") == 0
    matches = matches | untreated                # keep untreated participants
matches = matches | state.isna()                 # keep rows with no PDSTATE
return frame[matches.fillna(False)].copy()
\end{codeblock}

\begin{table}[h]
\centering
\begin{tabular}{@{}lrr@{}}
\toprule
\textbf{Operation} & \textbf{Rows} & \textbf{Removed} \\
\midrule
Part III raw                        & 39{,}310 & --- \\
After \code{select\_medication\_state} & 29{,}356 & 9{,}954 \\
\bottomrule
\end{tabular}
\end{table}

\begin{keypoint}
\textbf{Verified:} all 9{,}954 removed rows have \code{PDSTATE == 'ON'}. Not a
single row with \code{PDSTATE == 'OFF'}, with a missing \code{PDSTATE}, or from
an untreated participant was discarded. The filter does exactly what it claims.
\end{keypoint}

\textbf{Why this exists.} Levodopa suppresses tremor. Scoring an
on-medication examination would shift a patient's ratio toward PIGD and
manufacture transitions that are pharmacological rather than pathological. The
canonical Stebbins classification is defined on the OFF or untreated exam.

\textbf{Assumption introduced.} That an OFF exam reflects the underlying
untreated phenotype. In practice ``OFF'' in PPMI means a
practically-defined off state (typically overnight withdrawal), not a true
drug-free state.

\subsubsection*{Step 1.3 --- merge Part II with Part III \quad (\code{panel.py}:76)}

\begin{codeblock}
merged = part2[KEY + part2_items].merge(
    part3[KEY + part3_items + carry], on=KEY, how="outer"
)
\end{codeblock}

\begin{table}[h]
\centering
\begin{tabular}{@{}lrr@{}}
\toprule
\textbf{Operation} & \textbf{Rows} & \textbf{Duplicate keys} \\
\midrule
Part II (3 items) $\;\bowtie\;$ Part III OFF (13 items + 3 carried) & 37{,}897 & 11 \\
\bottomrule
\end{tabular}
\end{table}

\noindent
An \emph{outer} join is used, so a visit with only a questionnaire or only an
exam is retained rather than silently dropped; it will simply fail the
completeness test later and produce no label. Eleven duplicate keys survive the
medication filter --- these are visits with more than one qualifying OFF-state
record.

\subsubsection*{Step 1.4 --- \code{deduplicate\_visits()} \quad (\code{subtype.py}:187)}

\begin{codeblock}
present = tuple(c for c in cfg.all_items if c in frame.columns)
cleaned = clean_items(frame, present, cfg)
ordered = frame.assign(_n_missing=cleaned.isna().sum(axis=1))
ordered = ordered.sort_values(["PATNO", "EVENT_ID", "_n_missing"], kind="mergesort")
deduped = ordered.drop_duplicates(["PATNO", "EVENT_ID"], keep="first")
\end{codeblock}

\begin{table}[h]
\centering
\begin{tabular}{@{}lrr@{}}
\toprule
\textbf{Operation} & \textbf{Rows} & \textbf{Removed} \\
\midrule
After merge            & 37{,}897 & --- \\
After deduplication    & 37{,}886 & 11 \\
\bottomrule
\end{tabular}
\end{table}

\noindent
\textbf{The tie-break rule:} keep the row with the \emph{fewest missing items},
counted \emph{after} sentinel cleaning, and break remaining ties by first
occurrence using a stable sort (\code{kind="mergesort"}). The result does not
depend on the order pandas happened to read the file. This is unit-tested in
\code{tests/test\_subtype.py::TestDeduplication}.

\subsubsection*{Step 1.5 --- \code{compute\_subscores()} \quad (\code{subtype.py}:77)}

Covered in full in Section~\ref{sec:label}. Applied to all 37{,}886
patient-visits in the merged frame:

\begin{table}[h]
\centering
\begin{tabular}{@{}lr@{}}
\toprule
\code{tremor\_score} non-null & 27{,}046 \\
\code{pigd\_score} non-null   & 24{,}098 \\
\code{subtype} non-null       & \textbf{24{,}066} \\
\code{ratio\_defined} True    & 12{,}644 \\
\bottomrule
\end{tabular}
\end{table}

\noindent
Note that subscores are computed for the \emph{entire PPMI population}, not
only the analysis cohort. The restriction to sporadic PD happens at the next
step, via the join.

\subsubsection*{Step 1.6 --- the join onto the visit skeleton \quad (\code{panel.py}:69)}

\begin{codeblock}
panel = curated.merge(subscores, on=KEY, how="left", validate="one_to_one")
\end{codeblock}

\begin{table}[h]
\centering
\begin{tabular}{@{}lr@{}}
\toprule
Left side (cohort visits)  & 5{,}774 \\
Right side (all subscores) & 37{,}886 \\
\textbf{Result}            & \textbf{5{,}774 $\times$ 223} \\
\midrule
\quad of which labelled    & \textbf{4{,}441} \quad (76.9\%) \\
\bottomrule
\end{tabular}
\end{table}

\begin{keypoint}
\textbf{The direction of this join is a design decision worth defending.} It is
a \emph{left} join from the curated cut, so the row count is fixed by the visit
skeleton, not by the item files. A patient-visit exists in the panel if and only
if PPMI's curated cut says the visit happened. Item data can only \emph{fill in}
a label, never add or remove a row. The \code{validate="one\_to\_one"} argument
makes pandas raise if either side has duplicate keys --- an assertion that the
deduplication actually worked.
\end{keypoint}

\subsubsection*{Step 1.7 --- \code{\_add\_time\_axis()} \quad (\code{panel.py}:97)}

Covered in Section~\ref{sec:longitudinal}.

\subsection{D. Data shape after each step}

\begin{codeblock}
CURATED CUT                                  MDS-UPDRS PART II    MDS-UPDRS PART III
19,450 x 212                                 37,777 x 22          39,310 x 65
     |                                            |                    |
     | subgroup == "Sporadic PD"                  |                    | select_medication_state
     v                                            |                    v
5,789 x 212                                       |               29,356 x 65
     |                                            |                    |
     | drop EVENT_ID in {ST, PW}                  +--------+-----------+
     v                                                     | outer merge on (PATNO, EVENT_ID)
5,774 x 212                                                v
     |                                              37,897 x 21
     |                                                     | deduplicate_visits
     |                                                     v
     |                                              37,886 x 21
     |                                                     | compute_subscores
     |                                                     v
     |                                              24,066 labelled of 37,886
     |                                                     |
     +---------------------- left join --------------------+
                             v
                    5,774 x 223   <-- outputs/visit_panel.csv
                             |
                             | subtype.notna()
                             v
                    4,441 labelled visits  (76.9%)
\end{codeblock}

\subsection{E. Validation \quad (\code{validate.py})}
\label{sec:validation}

Two independent checks run before the panel is written.

\begin{table}[h]
\centering
\begin{tabular}{@{}llr@{}}
\toprule
\textbf{Reference} & \textbf{Classes} & \textbf{Agreement} \\
\midrule
Curated cut \code{td\_pigd}   & 2 (TD / non-TD)         & 4{,}431 / 4{,}431 = \textbf{100.00\%} \\
Archived \code{td\_pigd\_old} & 3 (TD / Indet. / PIGD)  & 1{,}595 / 1{,}597 = \textbf{99.87\%} \\
\bottomrule
\end{tabular}
\end{table}

\noindent
The two-class check confirms the arithmetic, the item lists, the medication
filter and the 1.15 threshold. It cannot test the Indeterminate/PIGD boundary,
because the curated cut collapses those classes. The three-class check, over
390 overlapping patients, is the only test of that boundary.

\begin{keypoint}
\textbf{A trap that was hit and fixed.} The archived \code{td\_pigd\_old} codes
\textbf{2 = PIGD and 3 = Indeterminate}, the reverse of this project's coding.
Compared without remapping, agreement reads 65.9\% and looks like a broken
classifier. \code{validate.remap\_archived\_labels()} (line 68) applies the
correction, which is driven by \code{config.yaml} lines 80--83 and unit-tested
in \code{tests/test\_validate.py}. The remap is verified against the archived
data dictionary, not guessed.
\end{keypoint}

\subsection{F. Research reasoning}

\textbf{What the code does.} Restricts to a homogeneous cohort, selects the
clinically correct examination, deduplicates deterministically, computes a
published phenotype from raw items, and proves the result against two
independent references.

\textbf{Why the researcher chose to do it.} The proposal requires a
\emph{three-class} label and a \emph{continuous ratio} at every visit. PPMI's
curated cut supplies neither: it has no tremor subscore, and its
\code{td\_pigd} is two-class. Recomputation was not optional.

\textbf{Assumptions introduced.} (i) Stebbins 2013 is the right phenotype;
(ii) an OFF exam represents the untreated state; (iii) sporadic PD is the
appropriate population; (iv) a visit with any missing item yields no label.

\textbf{What could go wrong.} Assumption (iv) is the one that bites: it discards
1{,}333 of 5{,}774 visits, and not at random. This is examined in question Q6 of
Section~\ref{sec:defence}.

% =============================================================================
\section{Target / label construction}
\label{sec:label}
% =============================================================================

\subsection{The formula, exactly as implemented}

\[
  \text{tremor\_score} \;=\; \frac{1}{11}\sum_{i \in \mathcal{T}} x_i,
  \qquad
  \text{pigd\_score} \;=\; \frac{1}{5}\sum_{j \in \mathcal{P}} x_j,
  \qquad
  \text{ratio} \;=\; \frac{\text{tremor\_score}}{\text{pigd\_score}}
\]

where $\mathcal{T}$ is the eleven tremor items and $\mathcal{P}$ the five
gait/postural items. Both are \textbf{means, not sums}.

\[
\text{subtype} =
\begin{cases}
\text{TD} & \text{ratio} \geq 1.15 \\
\text{TD} & \text{pigd\_score} = 0 \;\wedge\; \text{tremor\_score} > 0 \\
\text{Indeterminate} & 0.90 < \text{ratio} < 1.15 \\
\text{Indeterminate} & \text{tremor\_score} = 0 \;\wedge\; \text{pigd\_score} = 0 \\
\text{PIGD} & \text{ratio} \leq 0.90
\end{cases}
\]

Codes are \code{TD = 1}, \code{Indeterminate = 2}, \code{PIGD = 3}
(\code{config.yaml} lines 71--74).

\subsection{Step 1 --- cleaning \quad (\code{clean\_items}, \code{subtype.py}:39)}

\begin{codeblock}
lo, hi = cfg.valid_score_range                     # (0, 4)
block = frame[list(items)].apply(pd.to_numeric, errors="coerce")
return block.where((block >= lo) & (block <= hi))
\end{codeblock}

\begin{keypoint}
\textbf{PPMI stores ``Unable to Rate'' as the value 101, in the same numeric
field as a real 0--4 rating}, in all thirteen Part III items. Averaged as a
score, a single 101 among the five PIGD items lifts their mean to roughly
20, drives the ratio to almost zero, and relabels a tremor-dominant visit as
PIGD. \code{NP3PSTBL} alone carries 3{,}206 such codes among OFF exams.

The implementation excludes them by \emph{range}, not by an equality test
against 101, so a sentinel PPMI introduces later cannot silently enter a
subscore. Six unit tests in \code{tests/test\_subtype.py::TestUnrateableCodes}
pin this down.
\end{keypoint}

\subsection{Step 2 --- subscores with a missingness policy \quad (\code{\_subscore}, line 62)}

\begin{codeblock}
block = clean_items(frame, items, cfg)
n_missing = block.isna().sum(axis=1)
score = block.mean(axis=1, skipna=True)
return score.where(n_missing <= max_missing)       # max_missing = 0
\end{codeblock}

\noindent
With \code{max\_missing\_tremor\_items: 0} and
\code{max\_missing\_pigd\_items: 0}, \textbf{any} missing constituent nullifies
the subscore. This matches PPMI's own rule and is why the two-class validation
reaches exactly 100\%.

\subsection{Step 3 --- classification \quad (\code{classify}, line 121)}

\begin{codeblock}
both_zero      = observed & (t == 0) & (p == 0)
pigd_zero_only = observed & (p == 0) & (t > 0)
divisible      = observed & (p != 0)

ratio = t / p.where(p != 0)

result[divisible & (ratio >= cfg.td_threshold)]   = labels["TD"]
result[divisible & (ratio <= cfg.pigd_threshold)] = labels["PIGD"]
result[divisible & (ratio > pigd_thr) & (ratio < td_thr)] = labels["Indeterminate"]

result[pigd_zero_only] = labels["TD"]              # applied last: cannot be overwritten
result[both_zero]      = labels["Indeterminate"]
\end{codeblock}

\noindent
The degenerate cases are handled by explicit masks applied \emph{after} the
ratio cases, so they cannot be overwritten by a NaN comparison. A property test
(\code{test\_every\_observed\_visit\_is\_classified}) sweeps a grid of
$17 \times 17$ score pairs and asserts no observed pair ever yields NaN.

\subsection{Edge cases, and how each is resolved}

\begin{table}[h]
\centering
\small
\begin{tabular}{@{}lll@{}}
\toprule
\textbf{Situation} & \textbf{Naive result} & \textbf{Implemented result} \\
\midrule
$\text{pigd}=0$, $\text{tremor}>0$   & $\infty$ or crash & TD \\
$\text{tremor}=0$, $\text{pigd}=0$   & NaN ($0/0$)       & Indeterminate \\
$\text{tremor}=0$, $\text{pigd}>0$   & ratio $=0$        & PIGD (via $\leq 0.90$) \\
ratio exactly $1.15$                 & ambiguous         & TD (boundary closed) \\
ratio exactly $0.90$                 & ambiguous         & PIGD (boundary closed) \\
any item missing                     & biased mean       & no label \\
item $= 101$                         & mean $\approx 20$ & treated as missing \\
\bottomrule
\end{tabular}
\end{table}

\subsection{A fully worked example from the real data}

\textbf{PATNO 3010, visit BL.} Raw items as they appear on disk:

\begin{codeblock}
Part II:   NP2TRMR = 1     NP2WALK = 1     NP2FREZ = 0
Part III:  PDSTATE = NaN   PDTRTMNT = 0            <-- untreated participant
           NP3PTRMR = 1  NP3PTRML = 1  NP3KTRMR = 1  NP3KTRML = 1
           NP3RTARU = 0  NP3RTALU = 0  NP3RTARL = 0  NP3RTALL = 0
           NP3RTALJ = 0  NP3RTCON = 0
           NP3GAIT  = 0  NP3FRZGT = 0  NP3PSTBL = 0
\end{codeblock}

\textbf{Note first} that \code{PDSTATE} is missing and \code{PDTRTMNT = 0}.
This participant was not on PD medication at baseline, so there is no ON/OFF
distinction to make. \code{select\_medication\_state()} keeps the row through
its \code{untreated} clause --- exactly the case its
\code{test\_untreated\_participants\_are\_kept} test covers.

\textbf{Tremor score.} Eleven items, values
$1,1,1,1,1,0,0,0,0,0,0$:
\[
  \text{tremor\_score} = \frac{1+1+1+1+1+0+0+0+0+0+0}{11} = \frac{5}{11} = 0.4545
\]

\textbf{PIGD score.} Five items, values $1, 0, 0, 0, 0$:
\[
  \text{pigd\_score} = \frac{1+0+0+0+0}{5} = \frac{1}{5} = 0.2000
\]

\textbf{Ratio.}
\[
  \text{ratio} = \frac{0.4545}{0.2000} = 2.2727
\]

\textbf{Classification.} $2.2727 \geq 1.15 \Rightarrow$ \textbf{TD (code 1)}.

This matches \code{outputs/visit\_panel.csv} exactly: \code{tremor\_score
= 0.454545}, \code{pigd\_score = 0.2}, \code{ratio = 2.272727},
\code{subtype\_name = TD}.

\subsection{The ratio is not always defined}

Of 4{,}441 labelled visits, \textbf{725 (16.3\%)} have
$\text{pigd\_score} = 0$. Those visits carry a valid \emph{label} (the rule
sends them to TD) but no finite \emph{ratio}. \code{compute\_subscores} records
this distinction explicitly:

\begin{codeblock}
out["ratio"]         = tremor / pigd.replace(0.0, np.nan)   # NaN, not inf
out["ratio_defined"] = both_observed & (pigd != 0)
\end{codeblock}

\noindent
This matters because the proposed hierarchical model consumes the
\emph{continuous ratio} as a noisy observation of a latent phenotype. For one
visit in six there is nothing to consume. The flag exists so that population
can be counted and handled deliberately rather than dropped silently.

% =============================================================================
\section{Longitudinal panel construction}
\label{sec:longitudinal}
% =============================================================================

\subsection{The time axis \quad (\code{\_add\_time\_axis}, \code{panel.py}:97)}

\begin{codeblock}
out["visit_ts"] = parse_visit_date(out["visit_date"])          # "%m/%Y"
out = out.sort_values(["PATNO", "visit_ts", "YEAR"], kind="mergesort")
baseline = out.groupby("PATNO")["visit_ts"].transform("min")
out["years_from_baseline"] = (out["visit_ts"] - baseline).dt.days / 365.25
out["years_from_baseline"] = out["years_from_baseline"].fillna(out["YEAR"].astype(float))
out["visit_index"] = out.groupby("PATNO").cumcount()
out["n_visits"]    = out.groupby("PATNO")["PATNO"].transform("size")
\end{codeblock}

\begin{description}[leftmargin=2.4cm,style=nextline]
  \item[Visit identity] \code{(PATNO, EVENT\_ID)}. \code{EVENT\_ID} is a label
    such as \code{BL}, \code{V04}, \code{V06}, not a time.
  \item[Baseline] Each patient's \emph{own earliest} visit date
    (\code{transform("min")}), not a fixed calendar date and not necessarily
    \code{EVENT\_ID == "BL"}.
  \item[Elapsed time] $(\text{visit date} - \text{baseline date})/365.25$ days.
  \item[Fallback] If a date is unparseable, the nominal \code{YEAR} is used, so
    a missing \code{INFODT} does not drop a visit from the time axis.
  \item[Sorting] Stable mergesort, so ties resolve reproducibly.
\end{description}

\begin{keypoint}
\textbf{Why actual dates and not \code{YEAR}?} \code{YEAR} is derived from
\code{EVENT\_ID} --- it is a \emph{visit label}, not elapsed time. Patients miss
visits and attend late. In the panel, \code{V04} occurs anywhere from 0.83 to
over 1.4 years after baseline. Using \code{YEAR} would assign every
\code{BL}\arr\code{V04} pair an interval of exactly 1.0, which is false, and
would make the interval-stratified analysis in Section~\ref{sec:stats}
meaningless.

\textbf{Limitation, stated:} PPMI releases \code{INFODT} as month and year only,
so every date is anchored to the first of the month. Intervals carry roughly
$\pm$ two weeks of noise --- negligible at annual spacing.
\end{keypoint}

\subsection{What one patient's trajectory looks like}

\textbf{PATNO 3010}, taken verbatim from \code{outputs/visit\_panel.csv}:

\begin{codeblock}
EVENT_ID  YEAR visit_date  years_from_baseline  tremor  pigd    ratio  subtype
      BL     0    06/2011             0.000000  0.4545   0.2   2.2727       TD
     V04     1    06/2012             1.002053  0.9091   0.2   4.5455       TD
     V06     2    07/2013             2.083504  0.8182   0.2   4.0909       TD
     V08     3    09/2014             3.252567     NaN   NaN      NaN      ---   <- unlabelled
     V10     4    06/2015             4.000000  1.6364   0.2   8.1818       TD
     V12     5    11/2016             5.420945     NaN   NaN      NaN      ---   <- unlabelled
     V14     7    08/2018             7.167693  0.2727   0.6   0.4545     PIGD
     V15     8    07/2019             8.082136  1.4545   0.8   1.8182       TD
     V17    10    05/2021             9.916496     NaN   NaN      NaN      ---
     V18    11    06/2022            11.000684     NaN   NaN      NaN      ---
\end{codeblock}

\noindent
Observe: \code{V06} is at 2.083 years, not 2.000. \code{V17} is at 9.92 years
despite being labelled year 10. \code{V12} is at 5.42 years. The nominal year
and the elapsed time genuinely differ.

\subsection{Patients with too few usable visits}

\begin{table}[h]
\centering
\begin{tabular}{@{}lr@{}}
\toprule
Patients with $\geq 1$ labelled visit        & 1{,}080 \\
Patients with exactly 1 labelled visit       & \textbf{146} \\
Patients with $\geq 2$ labelled visits       & \textbf{934} \\
\bottomrule
\end{tabular}
\end{table}

\noindent
The 146 single-visit patients remain in \code{visit\_panel.csv} --- they are not
deleted --- but they cannot contribute a transition, so they contribute nothing
to Stage 2. \code{restrict\_to\_analysable()} (\code{panel.py}:120) applies the
\code{min\_visits\_per\_patient: 2} rule when producing
\code{outputs/baseline.csv}, which therefore has \textbf{934 rows}.

% =============================================================================
\section{Transition construction}
% =============================================================================

\subsection{\code{consecutive\_pairs()} \quad (\code{transitions.py}:16)}

\begin{codeblock}
df = panel[panel[label_col].notna()].sort_values(["PATNO", time_col])
grouped = df.groupby("PATNO", sort=False)

out = pd.DataFrame({
    "PATNO":      df["PATNO"],
    "from_state": df[label_col],
    "to_state":   grouped[label_col].shift(-1),
    "from_time":  df[time_col],
    "to_time":    grouped[time_col].shift(-1),
    "from_event": df["EVENT_ID"],
    "to_event":   grouped["EVENT_ID"].shift(-1),
})
out = out[out["to_state"].notna()].copy()
out["interval_years"] = out["to_time"] - out["from_time"]
\end{codeblock}

Two decisions are encoded here:

\begin{enumerate}
  \item \textbf{Unlabelled visits are filtered out \emph{before} pairing}
        (line 1). They therefore \emph{bridge} rather than truncate: a patient
        labelled at years 0 and 2 but not 1 yields one pair with a two-year
        interval, not a broken sequence.
  \item \textbf{\code{shift(-1)} is applied within \code{groupby("PATNO")}},
        so the last visit of one patient can never pair with the first visit of
        the next. \code{tests/test\_transitions.py::test\_pairs\_are\_within\_patient}
        asserts this.
\end{enumerate}

\subsection{The same patient, transformed}

\textbf{PATNO 3010} --- unlabelled visits removed, then paired:

\begin{codeblock}
from_event to_event  from_state  to_state  interval_years
        BL      V04          TD        TD        1.002053
       V04      V06          TD        TD        1.081451
       V06      V10          TD        TD        1.916496   <- V08 was unlabelled: BRIDGED
       V10      V14          TD      PIGD        3.167693   <- V12 unlabelled: 3.17y interval
       V14      V15        PIGD        TD        0.914442   <- an immediate REVERSION
\end{codeblock}

\noindent
This single patient demonstrates three mechanisms at once: gap bridging
(\code{V06}\arr\code{V10}), a long-interval transition
(\code{V10}\arr\code{V14}, 3.17 years, not 3.0), and a state change that
reverses at the very next visit (\code{V14}\arr\code{V15}). The last is exactly
what the reversion analysis in Section~\ref{sec:stats} counts.

\subsection{Arithmetic check on the pair count}

\begin{keypoint}
A patient with $n_i$ labelled visits contributes $n_i - 1$ pairs. Summed over
the 934 patients with $n_i \geq 2$:
\[
  \sum_i (n_i - 1) = 3{,}361
\]
which is exactly the number of rows \code{consecutive\_pairs()} returns.
\textbf{Verified independently.} No pair is lost or double-counted.
\end{keypoint}

\subsection{From counts to probabilities \quad (\code{transition\_matrix}, line 47)}

\begin{codeblock}
counts = pd.crosstab(pairs["from_state"], pairs["to_state"])
counts = counts.reindex(index=states, columns=states, fill_value=0)
totals = counts.sum(axis=1)
probs  = counts.div(totals.where(totals > 0), axis=0)
\end{codeblock}

\textbf{Raw counts, before normalisation} (total $= 3{,}361$):

\begin{table}[h]
\centering
\begin{tabular}{@{}lrrr|r@{}}
\toprule
& \textbf{$\rightarrow$ TD} & \textbf{$\rightarrow$ Ind} & \textbf{$\rightarrow$ PIGD} & \textbf{Row total} \\
\midrule
\textbf{TD}            & 1827 & 181 & 253 & 2261 \\
\textbf{Indeterminate} &  147 &  64 & 123 &  334 \\
\textbf{PIGD}          &  141 &  87 & 538 &  766 \\
\bottomrule
\end{tabular}
\end{table}

\textbf{After dividing each row by its total:}

\begin{table}[h]
\centering
\begin{tabular}{@{}lrrr@{}}
\toprule
& \textbf{$\rightarrow$ TD} & \textbf{$\rightarrow$ Ind} & \textbf{$\rightarrow$ PIGD} \\
\midrule
\textbf{TD}            & 0.808 & 0.080 & 0.112 \\
\textbf{Indeterminate} & 0.440 & 0.192 & 0.368 \\
\textbf{PIGD}          & 0.184 & 0.114 & 0.702 \\
\bottomrule
\end{tabular}
\end{table}

\noindent
Worked entry: $181 / 2261 = 0.0800$. The \code{reindex} guarantees all three
states appear even if one were unobserved, and \code{totals.where(totals > 0)}
prevents a division by zero for an empty row.

% =============================================================================
\section{Statistical analysis}
\label{sec:stats}
% =============================================================================

\subsection{Wilson score intervals \quad (\code{transition\_counts\_with\_ci}, line 67)}

\textbf{Question answered.} How precise is each transition probability?

\textbf{Formula, as implemented:}
\[
  \hat p = \frac{k}{n}, \qquad
  \text{centre} = \frac{\hat p + \frac{z^2}{2n}}{1 + \frac{z^2}{n}}, \qquad
  \text{half-width} = \frac{z\sqrt{\frac{\hat p(1-\hat p)}{n} + \frac{z^2}{4n^2}}}{1 + \frac{z^2}{n}}
\]
with $z = \Phi^{-1}(1 - \alpha/2)$, $\alpha = 0.05$.

\textbf{Why not the normal approximation?} Row counts range from 334 to
2{,}261, and several probabilities sit near 0 or 1 (\code{TD}\arr\code{Ind} is
0.080). The Wald interval $\hat p \pm z\sqrt{\hat p(1-\hat p)/n}$ produces
bounds outside $[0,1]$ in exactly those cells. A unit test
(\code{test\_wilson\_interval\_stays\_in\_unit\_range}) asserts the
implementation never does.

\textbf{Interpretation.} \code{TD}\arr\code{TD} is $0.808$, CI
$[0.791, 0.824]$. \textbf{Incorrect interpretation:} treating these as
independent --- the three probabilities in a row are constrained to sum to 1,
so their errors are negatively correlated.

\subsection{Interval stratification \quad (line 106)}

\textbf{Question answered.} Does a transition probability mean anything without
saying over what period?

Bins: $[0, 1.5)$, $[1.5, 2.5)$, $[2.5, 4.0)$, $[4.0, \infty)$ years.

\begin{table}[h]
\centering
\begin{tabular}{@{}lrr@{}}
\toprule
\textbf{Interval (years)} & \textbf{Pairs} & \textbf{TD $\rightarrow$ PIGD} \\
\midrule
$\leq 1.5$    & 2{,}882 & 0.100 \\
$1.5$--$2.5$  &    302 & 0.168 \\
$2.5$--$4.0$  &    111 & 0.153 \\
$> 4.0$       &     66 & 0.289 \\
\bottomrule
\end{tabular}
\end{table}

\noindent
\textbf{Conclusion this forces.} The pooled matrix is not a well-defined
quantity --- it averages one-year and multi-year transitions in whatever
proportion attendance happened to produce. It describes the study design as
much as the disease.

\subsection{Reversion analysis \quad (\code{reversion\_after\_change}, line 198)}

\textbf{Question answered.} Of the state changes we observe, how many are real?

\begin{codeblock}
frame = pd.DataFrame({
    "s0": df[label_col],
    "s1": g[label_col].shift(-1),
    "s2": g[label_col].shift(-2),
    ...})
frame = frame[frame["s2"].notna() & (frame["s0"] != frame["s1"])]
frame["reverted"] = frame["s2"] == frame["s0"]
\end{codeblock}

\noindent
It takes \emph{triples}: a change $s_0 \neq s_1$ that has a third labelled
visit $s_2$, and asks whether $s_2 = s_0$.

\begin{table}[h]
\centering
\begin{tabular}{@{}lr@{}}
\toprule
Qualifying changes (denominator) & 673 \\
Of which reverted                & 237 \\
\textbf{Reversion rate}          & \textbf{0.3522} \\
Patients contributing            & 342 \\
\bottomrule
\end{tabular}
\end{table}

\textbf{Assumption.} That a change followed by a return is more likely
measurement noise than two genuine changes. \textbf{Incorrect
interpretation:} concluding the labels are worthless. Section~4.3 of the
Weeks 3--4 report shows the movement is significantly \emph{asymmetric}
(TD\arr PIGD 253 versus PIGD\arr TD 141, $p = 1.8\times 10^{-8}$), so a real
directional drift exists beneath the churn.

\textbf{Limitation.} Requires three labelled visits, so the estimate comes
from the better-attended subset --- 342 of 934 patients.

\subsection{Kaplan--Meier \quad (\code{first\_entry} + \code{kaplan\_meier}, lines 246, 299)}

\textbf{Question answered.} How long until a patient first becomes PIGD?

\textbf{Formula, as implemented:}
\[
  \hat S(t) = \prod_{t_i \leq t}\left(1 - \frac{d_i}{n_i}\right),
  \qquad
  \widehat{\mathrm{SE}}(\hat S(t)) = \hat S(t)\sqrt{\sum_{t_i \leq t}\frac{d_i}{n_i(n_i - d_i)}}
\]
(Greenwood's formula). Cumulative incidence is plotted as $1 - \hat S(t)$.

\textbf{Risk set construction.} \code{first\_entry()} returns every patient
with a \code{converted} flag and a time; patients already PIGD at their first
visit are flagged \code{starts\_in\_state} and excluded.

\begin{table}[h]
\centering
\begin{tabular}{@{}lrr@{}}
\toprule
& \textbf{Naive (1 visit)} & \textbf{Sustained (2 visits)} \\
\midrule
Patients                       & 1{,}080 & 1{,}080 \\
Excluded (PIGD at baseline)    & \phantom{0,}195 & \phantom{0,}195 \\
\textbf{At risk}               & \phantom{0,}\textbf{885} & \phantom{0,}\textbf{885} \\
Converted                      & \phantom{0,}263 & \phantom{0,}127 \\
KM cumulative incidence, yr 5  & \textbf{35.25\%} & \textbf{15.33\%} \\
$n$ at risk at year 5          & \phantom{0,}180 & \phantom{0,}211 \\
\bottomrule
\end{tabular}
\end{table}

\textbf{Why Kaplan--Meier and not a percentage?} Follow-up ranges from one
visit to fourteen years. A raw proportion conflates ``did not convert'' with
``was not observed long enough''. KM censors at the last labelled visit.

\textbf{The \code{min\_consecutive} parameter.} With
\code{min\_consecutive=2}, an entry counts only if the patient is still PIGD at
the next labelled visit. Requiring this removes \textbf{136 of 263} apparent
conversions --- 51.7\%.

\textbf{Assumption.} Non-informative censoring: that dropping out is unrelated
to subtype. \textbf{This is questionable here} and is flagged as a limitation.

\subsection{Pathway analysis \quad (\code{state\_sequences} + \code{collapse\_runs}, lines 177, 185)}

\begin{codeblock}
def collapse_runs(sequence):
    out = []
    for s in sequence:
        if not out or out[-1] != s:
            out.append(s)
    return tuple(out)
\end{codeblock}

\noindent
\code{(TD, TD, PIGD, PIGD)} becomes \code{(TD, PIGD)}. \textbf{Why:} without
collapsing, pathway frequencies would be dominated by how often a patient
attended rather than by what happened to them.

\begin{table}[h]
\centering
\begin{tabular}{@{}lrr@{}}
\toprule
Patients with $\geq 2$ labels & 934 & \\
Distinct collapsed pathways   & 86 & \\
No recorded change of state   & 476 & 51.0\% \\
\bottomrule
\end{tabular}
\end{table}

% =============================================================================
\section{Trace of every major result}
% =============================================================================

Each trace follows the form \textsc{result} \arr\ \textsc{script} \arr\
\textsc{function} \arr\ \textsc{dataframe} \arr\ \textsc{operation} \arr\
\textsc{number}.

\subsection*{19{,}450 visits}
\begin{codeblock}
loaders.load_curated_cut(cfg)
  -> pd.read_excel("data/clinical/PPMI_Curated_Data_Cut_Public_20260511.xlsx",
                   sheet_name="20260511")
  -> df.shape = (19450, 212)
\end{codeblock}

\subsection*{5{,}774 cohort visits}
\begin{codeblock}
01_build_panel.py -> build_visit_panel() -> panel.filter_cohort()
  frame[frame["subgroup"].isin({"Sporadic PD"})]           19,450 -> 5,789
  frame[~frame["EVENT_ID"].isin({ST, PW, U01, U02})]        5,789 -> 5,774
  frame[frame["YEAR"] <= 15]                                5,774 -> 5,774
  SAVED: outputs/visit_panel.csv  (5,774 rows x 223 cols)
\end{codeblock}

\subsection*{29{,}356 OFF-state exam rows}
\begin{codeblock}
panel.build_subscores() -> subtype.select_medication_state()
  keep if PDSTATE == "OFF"  OR  PDTRTMNT == 0  OR  PDSTATE is NaN
  39,310 -> 29,356    (9,954 removed, all PDSTATE == "ON")
\end{codeblock}

\subsection*{4{,}441 labelled visits}
\begin{codeblock}
panel.build_visit_panel()
  curated(5,774) .merge(subscores, how="left", validate="one_to_one")  -> 5,774
  panel["subtype"].notna().sum()                                       -> 4,441
  = 76.9% of cohort visits
  SAVED: outputs/visit_panel.csv, column `subtype`
  REPORTED: outputs/panel_summary.txt
\end{codeblock}

\subsection*{934 patients / 146 excluded}
\begin{codeblock}
labelled = panel[panel.subtype.notna()]
per_patient = labelled.groupby("PATNO").size()
  (per_patient == 1).sum()  -> 146     cannot form a transition
  (per_patient >= 2).sum()  -> 934
  SAVED: outputs/baseline.csv  (934 rows, via panel.restrict_to_analysable)
\end{codeblock}

\subsection*{3{,}361 transition pairs}
\begin{codeblock}
02_transition_analysis.py -> transitions.consecutive_pairs(panel)
  df = panel[panel["subtype"].notna()].sort_values(["PATNO","years_from_baseline"])
  to_state = df.groupby("PATNO")["subtype"].shift(-1)
  out = out[out["to_state"].notna()]
  len(out) -> 3,361      CHECK: sum(n_i - 1) over 934 patients = 3,361
\end{codeblock}

\subsection*{0.808 (TD $\rightarrow$ TD)}
\begin{codeblock}
transitions.transition_matrix(pairs, states=[1,2,3])
  counts = pd.crosstab(pairs["from_state"], pairs["to_state"])
  counts.loc[TD, TD]  = 1827
  counts.loc[TD].sum()= 2261
  1827 / 2261 = 0.8080
  SAVED: outputs/transitions/transition_matrix_probs.csv
\end{codeblock}

\subsection*{35.2\% reversion}
\begin{codeblock}
transitions.reversion_after_change(panel)
  triples (s0, s1, s2) where s2 is not NaN AND s0 != s1     -> 673 rows
  reverted = (s2 == s0)                                     -> 237 True
  237 / 673 = 0.3522
  SAVED: outputs/transitions/reversion_summary.csv
\end{codeblock}

\subsection*{35.3\% vs 15.3\% conversion at year 5}
\begin{codeblock}
transitions.first_entry(panel, target_state=3, min_consecutive=k)
  at_risk = fe[~fe["starts_in_state"]]                      -> 885 patients
  km = transitions.kaplan_meier(at_risk.time, at_risk.converted)
  1 - km[km.time <= 5.0].survival.iloc[-1]
      k=1 -> 0.3525   (n at risk 180)
      k=2 -> 0.1533   (n at risk 211)
  SAVED: outputs/transitions/conversion_km_{naive,sustained}.csv
\end{codeblock}

\subsection*{100.00\% and 99.87\% validation}
\begin{codeblock}
01_build_panel.py -> validate.against_curated_cut(panel, cfg)
  ours_binary = collapse_to_binary(panel["subtype"])   # TD->1, else 2
  compare vs panel["td_pigd"]        -> 4,431 compared, 4,431 agree
01_build_panel.py -> validate.against_archived_three_class(panel, cfg)
  ref = archived[["PATNO","EVENT_ID","td_pigd_old"]]
  ref = remap_archived_labels(ref)   # 2<->3 swap, per archived dictionary
  inner join, compare               -> 1,597 compared, 1,595 agree
  SAVED: outputs/validation.txt
\end{codeblock}

% =============================================================================
\section{Complete data funnel}
% =============================================================================

\begin{codeblock}
PPMI CURATED DATA CUT                                        19,450 visits (4,788 patients)
     |
     |  subgroup == "Sporadic PD"        -- other cohorts (Hyposmia, RBD, HC,
     |                                      LRRK2, GBA, SWEDD...) not selected
     v
                                                              5,789 visits (1,080 patients)
     |  drop EVENT_ID in {ST, PW}        -- unscheduled / premature-withdrawal visits
     v
COHORT VISIT SKELETON                                         5,774 visits (1,080 patients)
     |
     |  ---------------- meanwhile, on the item files ----------------
     |
     |   MDS-UPDRS Part II   37,777 rows   (0 duplicate visit keys)
     |   MDS-UPDRS Part III  39,310 rows   (6,907 duplicate visit keys)
     |         |
     |         |  select_medication_state: keep OFF / untreated
     |         v
     |   Part III OFF        29,356 rows   (9,954 ON-state rows removed)
     |         |
     |         |  outer merge on (PATNO, EVENT_ID)
     |         v
     |                       37,897 rows   (11 duplicate keys remain)
     |         |
     |         |  deduplicate_visits: fewest missing items wins
     |         v
     |                       37,886 patient-visits
     |         |
     |         |  clean_items (drop 101 / out-of-range) -> subscores -> classify
     |         v
     |                       24,066 labelled  (of 37,886, all PPMI patients)
     |         |
     +---------+  LEFT JOIN onto the cohort skeleton
     v
LABELLED PANEL                                     5,774 rows x 223 cols
                                                   outputs/visit_panel.csv
     |
     |  subtype.notna()      -- 1,333 visits have no complete OFF-state item set
     v
                                                   4,441 labelled visits (76.9%)
     |
     |  group by PATNO, sort by years_from_baseline
     v
                                                   1,080 patients
     |  keep patients with >= 2 labelled visits    -- 146 single-visit patients
     v                                                cannot form a transition
                                                   934 patients
     |
     |  consecutive_pairs: shift(-1) within patient
     v
TRANSITION PAIRS                                   3,361 pairs
     |
     +--> transition_matrix        -> 3x3 counts and probabilities
     +--> transition_..._with_ci   -> Wilson 95% intervals
     +--> ..._by_interval          -> 4 interval strata (2,882 / 302 / 111 / 66)
     +--> reversion_after_change   -> 673 changes, 237 reverted = 35.2%
     +--> first_entry + kaplan_meier -> 885 at risk, 263 / 127 converted
     +--> state_sequences          -> 86 pathways over 934 patients, 476 unchanged
     v
FINAL RESULTS                                      outputs/transitions/
                                                   9 CSV tables, 6 figures, report.txt
\end{codeblock}

% =============================================================================
\section{Code $\leftrightarrow$ science map}
% =============================================================================

All modules live in \code{src/ppmi\_subtype/}.

\begin{longtable}{@{}p{2.3cm} p{2.2cm} p{3.7cm} p{2.9cm} p{2.4cm}@{}}
\toprule
\textbf{Concept} & \textbf{Module} & \textbf{Function} & \textbf{Operation} & \textbf{Output} \\
\midrule
\endfirsthead
\toprule
\textbf{Concept} & \textbf{Module} & \textbf{Function} & \textbf{Operation} & \textbf{Output} \\
\midrule
\endhead
Study population & \code{panel} & \code{filter\_cohort} & subgroup + event filter & 5{,}774 visits \\
Correct exam state & \code{subtype} & \code{select\_}\newline\code{medication\_state} & keep OFF or untreated & 29{,}356 rows \\
One record per visit & \code{subtype} & \code{deduplicate\_visits} & fewest-missing tie-break & 37{,}886 rows \\
Invalid measurements & \code{subtype} & \code{clean\_items} & drop outside 0--4 & cleaned items \\
Motor phenotype & \code{subtype} & \code{compute\_subscores} & two means + ratio & tremor, pigd, ratio \\
Clinical subtype & \code{subtype} & \code{classify} & 1.15 / 0.90 thresholds & TD / Ind / PIGD \\
Label correctness & \code{validate} & \code{against\_curated\_cut}, \code{against\_archived\_}\newline\code{three\_class} & agreement & 100\% and 99.87\% \\
Disease timeline & \code{panel} & \code{\_add\_time\_axis} & date $-$ patient min & \code{years\_from\_}\newline\code{baseline} \\
Patient trajectory & \code{transitions} & \code{state\_sequences} & groupby + tuple & state sequence \\
Observed transition & \code{transitions} & \code{consecutive\_pairs} & groupby + shift($-1$) & 3{,}361 pairs \\
Progression rates & \code{transitions} & \code{transition\_matrix} & crosstab + row-normalise & $3\times3$ matrix \\
Estimate precision & \code{transitions} & \code{transition\_counts\_}\newline\code{with\_ci} & Wilson score & 95\% CIs \\
Time-dependence & \code{transitions} & \code{transition\_matrix\_}\newline\code{by\_interval} & cut + per-bin matrix & 4 matrices \\
Measurement noise & \code{transitions} & \code{reversion\_}\newline\code{after\_change} & triples $s_0,s_1,s_2$ & 35.2\% \\
Time to event & \code{transitions} & \code{first\_entry}, \code{kaplan\_meier} & survival + censoring & KM curves \\
Common courses & \code{transitions} & \code{collapse\_runs} & drop repeated states & 86 pathways \\
Figures & \code{plots} & six plotting functions & matplotlib & 6 PNG files \\
\bottomrule
\end{longtable}

% =============================================================================
\section{Professor defence questions}
\label{sec:defence}
% =============================================================================

\subsection*{Q1. Why did you recompute the label instead of using PPMI's \code{td\_pigd}?}
\textbf{Answer.} PPMI's curated \code{td\_pigd} is coded two-class (1 = TD,
2 = non-TD), collapsing Indeterminate away, and the cut ships \code{pigd} with
no matching tremor subscore --- the ratio has a denominator and no numerator.
The proposal requires all three classes and the continuous ratio.

\textbf{Evidence in code.} \code{docs/DATA\_REQUIREMENTS.md}; the curated
dictionary decode for \code{td\_pigd}.

\textbf{Follow-up: how do you know your recomputation is right?} Two
independent checks. Collapsed to binary it agrees with PPMI's own
\code{td\_pigd} on 4{,}431 of 4{,}431 visits. Against the archived three-class
\code{td\_pigd\_old} it agrees on 1{,}595 of 1{,}597.

\subsection*{Q2. Why OFF-medication visits?}
\textbf{Answer.} Levodopa suppresses tremor. Scoring an ON exam shifts the
ratio toward PIGD and manufactures transitions that are pharmacological, not
pathological. The Stebbins definition is specified on the OFF or untreated exam.

\textbf{Evidence.} \code{subtype.select\_medication\_state}; verified that all
9{,}954 removed rows are \code{PDSTATE == 'ON'}.

\textbf{Follow-up: don't you lose untreated patients?} No --- the filter has an
explicit \code{PDTRTMNT == 0} clause, unit-tested by
\code{test\_untreated\_participants\_are\_kept}. At baseline 99.7\% of the
cohort is untreated and every one is retained.

\subsection*{Q3. Why actual dates rather than visit numbers?}
\textbf{Answer.} \code{YEAR} is derived from \code{EVENT\_ID} --- a visit label,
not elapsed time. In the real panel \code{V06} lands at 2.08 years and
\code{V17} at 9.92. Using \code{YEAR} would force every interval to a whole
number and destroy the interval-stratified analysis.

\textbf{Evidence.} \code{panel.\_add\_time\_axis}; PATNO 3010's trajectory.

\textbf{Follow-up: how precise are the dates?} Month resolution only ---
roughly $\pm$ two weeks, negligible at annual spacing. Stated as a limitation.

\subsection*{Q4. Why do unlabelled visits bridge instead of breaking the sequence?}
\textbf{Answer.} Because a missed \emph{assessment} is not a missed
\emph{patient}. PATNO 3010 has no label at \code{V08}; treating that as an
endpoint would discard six years of subsequent follow-up. Instead
\code{V06}\arr\code{V10} becomes one pair with a 1.92-year interval, and the
interval is carried so the longer gap is visible to the analysis.

\textbf{Evidence.} \code{transitions.consecutive\_pairs} line 1, and
\code{test\_unlabelled\_visits\_are\_skipped\_not\_terminal}.

\textbf{Follow-up: doesn't that assume nothing happened in between?} It assumes
nothing \emph{observed} happened. That is precisely why intervals are
stratified rather than pooled --- a four-year pair is not treated as equivalent
to a one-year pair.

\subsection*{Q5. Why not impute the missing subtype?}
\textbf{Answer.} \textbf{The implementation does not impute. The rationale is
not explicitly documented in the repository.} What \emph{is} documented and
tested is the sensitivity check: relaxing the rule to tolerate one missing item
recovers only 44 visits and changes zero existing labels.

\textbf{Interpretation (labelled as such).} Imputing the subtype would mean
imputing the outcome variable, which risks manufacturing the transitions the
study is trying to measure.

\subsection*{Q6. Why are there fewer labelled visits than cohort visits?}
\textbf{Answer.} 1{,}333 of 5{,}774 visits (23.1\%) have no complete
OFF-medication item set. Of these, 722 have no OFF exam recorded at all, 595
have an exam that could not be fully scored, and 16 lack Part II.

\textbf{Follow-up: could this introduce selection bias?} \textbf{Yes, and its
direction is known.} Among the 1{,}253 unscorable OFF exams, \code{NP3PSTBL}
--- item 3.12, postural stability --- is the missing item in 1{,}243 of them.
The pull test is often unsafe in a patient with poor balance. Those exams come
from more advanced patients: median Hoehn \& Yahr 3.0 against 2.0
($p = 1.4\times10^{-17}$). Postural stability is a \emph{PIGD} item, so the
visits lost are disproportionately those most likely to be PIGD. \textbf{The
reported drift toward PIGD is therefore a lower bound.}

\subsection*{Q7. Why Wilson intervals?}
\textbf{Answer.} Because one estimator has to serve both the pooled matrix and
the interval-stratified table, and the normal (Wald) approximation breaks down
in the second. In the pooled $3\times3$ matrix Wald would in fact be adequate
--- checked cell by cell, it never leaves $[0,1]$ and agrees with Wilson to
within 0.004. But the stratified table has row counts as low as 9, and there
Wald returns a \emph{negative} lower bound in five cells, for example
Indeterminate\arr Indeterminate at $2.5$--$4.0$ years ($k=2$, $n=16$):
Wald $[-0.037,\, 0.287]$ against Wilson $[0.035,\, 0.360]$.

Even where Wald stays inside $[0,1]$ it is wrong at small $n$, because it
forces a symmetric interval around $\hat p$ when the true uncertainty is
skewed. Switching estimator between tables would be worse than using the one
that is valid for both.

\textbf{Evidence.} \code{transition\_counts\_with\_ci};
\code{test\_wilson\_interval\_stays\_in\_unit\_range}.

\subsection*{Q8. Why Kaplan--Meier instead of a conversion percentage?}
\textbf{Answer.} Follow-up ranges from one visit to fourteen years. A raw
percentage cannot distinguish ``did not convert'' from ``was not observed long
enough''. KM censors at the last labelled visit.

\textbf{Follow-up: is censoring non-informative here?} Probably not, and we say
so. Dropout is heavy and, per Q6, scorability is tied to severity. The curves
are truncated at year 12 where the risk set falls below 25.

\subsection*{Q9. Why a transition matrix at all?}
\textbf{Answer.} It is the empirical object the Week 5--7 model must reproduce.
Before fitting any model, we need to know what transitions look like: which
states persist, which are transient, and whether movement is directional.

\textbf{Follow-up: is the process Markov?} \textbf{Not established, and there
is evidence against it.} Iterating the pooled matrix gives a stationary
distribution of TD 0.560 / Ind 0.103 / PIGD 0.338, whereas observed year-11
prevalence is 0.391 / 0.120 / 0.489. The gap is reported rather than hidden.

\subsection*{Q10. Does a TD $\rightarrow$ PIGD transition prove biological progression?}
\textbf{Answer.} No, and the analysis is built to show why not. 35.2\% of all
observed changes reverse at the very next visit. A single observed transition
is weak evidence about one patient.

\textbf{Follow-up: then is any of it real?} Yes. The movement is significantly
asymmetric: TD\arr PIGD occurs 253 times against 141 in reverse,
$p = 1.8\times10^{-8}$. Of 932 cross-state moves, 80.5\% cancels out and 19.5\%
is net drift toward PIGD. The correct reading is a genuine drift observed
through a noisy measurement --- which is the case \emph{for} a latent-state
model, not against it.

\subsection*{Q11. Why patient-level rather than visit-level reasoning?}
\textbf{Answer.} Transitions are only defined within a patient.
\code{consecutive\_pairs} does \code{shift(-1)} inside
\code{groupby("PATNO")}, so one patient's last visit can never pair with the
next patient's first. Patients with a single labelled visit (146) contribute
nothing.

\textbf{Follow-up: are patients weighted equally?} \textbf{No.} A patient with
14 visits contributes 13 pairs; one with 2 contributes 1. The pooled matrix is
therefore visit-weighted, not patient-weighted, and better-attended patients
carry more influence. \textbf{This is not currently corrected for} and is a
fair criticism.

% =============================================================================
\section{What we have actually achieved}
% =============================================================================

\begin{enumerate}
  \item \textbf{A validated phenotype.} A three-class TD/PIGD label plus the
        continuous ratio, computed from raw MDS-UPDRS items for 4{,}441 visits,
        agreeing with PPMI's own derivations at 100\% (2-class) and 99.87\%
        (3-class).
  \item \textbf{A longitudinal panel.} 1{,}080 patients, 5{,}774 visits, real
        elapsed-time axis, out to 14 years --- \code{outputs/visit\_panel.csv}.
  \item \textbf{An empirical characterisation of transitions.} 3{,}361 pairs;
        matrices with Wilson intervals; interval stratification; reversion;
        time-to-conversion under two definitions; pathway analysis.
  \item \textbf{Four findings that constrain the model.} Indeterminate is a
        crossing ($p_{\text{stay}} = 0.192$), PIGD is not absorbing (29.8\%
        leave), 35.2\% of changes reverse immediately, and transition
        probabilities scale with the inter-visit interval.
  \item \textbf{A named, directional bias.} Label loss is concentrated in
        patients who cannot complete the pull test, making the reported drift a
        lower bound.
  \item \textbf{Reproducibility.} Two commands regenerate every number;
        74 unit tests; all parameters in one config file.
\end{enumerate}

\begin{keypoint}
\textbf{What has \emph{not} been done.} No model has been fitted. No parameter
has been estimated. No prediction has been made. Everything above is
descriptive characterisation of the observed data.
\end{keypoint}

% =============================================================================
\section{What comes next}
% =============================================================================

Weeks 5--7 in the proposal: a baseline probabilistic model. The analysis in
this document changes what that baseline should be.

\begin{description}[leftmargin=1.6cm,style=nextline]
  \item[Continuous-time Markov chain, not discrete] Because TD\arr PIGD ranges
    from 0.100 to 0.289 depending only on the gap between visits, a chain
    indexed by ``visit number'' would encode PPMI's attendance schedule.
    Estimating a generator matrix $Q$ and using $P(t) = e^{Qt}$ handles
    arbitrary intervals.
  \item[Decide the ratio question] 16.3\% of labelled visits have
    $\text{pigd\_score} = 0$ and therefore no finite ratio. The recommended
    route is modelling the two subscores jointly and deriving the ratio's
    posterior, which keeps every visit.
  \item[Address the measurement noise explicitly] 35.2\% immediate reversion is
    the empirical justification for the proposal's
    $\text{observed} = \text{latent} + \text{noise}$ structure.
  \item[Quantify the missingness bias] Impute \code{NP3PSTBL}, or model the
    label as missing-not-at-random, to put a size on the lower bound in Q6.
  \item[Then the hierarchical Bayesian model] Population-level progression plus
    subject-level random effects, as in the proposal, followed by the
    $\sigma_i$-as-phenotype analysis in its second half.
\end{description}

\vspace{1em}
\noindent\textbf{Reference.} Stebbins GT, Goetz CG, Burn DJ, Jankovic J, Khoo TK,
Tilley BC. How to identify tremor dominant and postural instability/gait
difficulty groups with the Movement Disorder Society Unified Parkinson's Disease
Rating Scale. \emph{Movement Disorders} 2013;28(5):668--70.

\end{document}
